\documentclass[11pt,letterpaper]{article}

%% ── Packages ────────────────────────────────────────────────────────────────
\usepackage[margin=1in]{geometry}
\usepackage{amsmath,amssymb,bm}
\usepackage{graphicx}
\usepackage{booktabs}
\usepackage{hyperref}
\usepackage{caption}
\usepackage{subcaption}
\usepackage{xcolor}
\usepackage{enumitem}
\usepackage{microtype}
\usepackage{parskip}

\hypersetup{
  colorlinks=true, linkcolor=blue, citecolor=blue, urlcolor=blue
}

%% ── Macros ──────────────────────────────────────────────────────────────────
\newcommand{\R}{\mathbb{R}}
\newcommand{\W}{\mathbf{W}}
\newcommand{\psi}{\boldsymbol{\psi}}
\renewcommand{\a}{\mathbf{a}}

%% ── Title ───────────────────────────────────────────────────────────────────
\title{\textbf{ECE 66200: Pattern Recognition — Spring 2026}\\[4pt]
       \Large Project 1, Part 1: PCA View of CNN Filters (DCF in AlexNet)}
\author{[Your Name] \\ Purdue University}
\date{\today}

\begin{document}
\maketitle
\tableofcontents
\newpage

% ═══════════════════════════════════════════════════════════════════════════════
\section{Problem Formulation}
\label{sec:formulation}
% ═══════════════════════════════════════════════════════════════════════════════

\subsection{Goal}

The central question of Part 1 is: \emph{do CNN convolutional filters occupy a
low-dimensional subspace, and does truncating that subspace behave like
Principal Component Analysis (PCA)?}

We answer this by (i) training an AlexNet on MNIST, (ii) applying PCA to the
learned filter tensors, (iii) replacing the convolution layers with
Decomposed Convolutional Filter (DCF) layers parameterised by
$K$ filter atoms, and (iv) measuring how test accuracy and PSNR vary with $K$.

\subsection{Notation}

Let $\W \in \R^{c' \times c \times l \times l}$ denote the weight tensor of a
convolutional layer with $c'$ output channels, $c$ input channels, and
$l \times l$ spatial kernels.  Reshaping to
$\W' \in \R^{c'c \times l^2}$ reveals the \emph{filter matrix}, whose rows
are individual $l \times l$ filters laid flat.

% ═══════════════════════════════════════════════════════════════════════════════
\section{Mathematical Description of DCF}
\label{sec:dcf_math}
% ═══════════════════════════════════════════════════════════════════════════════

\subsection{Decomposed Convolutional Filters (DCF)}

The DCF representation \cite{qiu2018dcfnet} expresses each convolutional
filter as a \emph{linear combination of $K$ shared atoms}:

\begin{equation}
  \W[i,j] \;=\; \sum_{k=1}^{K} a_k[i,j]\,\psi_k,
  \qquad
  \psi_k \in \R^{l \times l},\quad
  a_k[i,j] \in \R,
  \label{eq:dcf}
\end{equation}

where $\psi_k$ are the \textbf{filter atoms} (basis elements shared across
all input/output channel pairs) and $a_k \in \R^{c' \times c}$ is the
\textbf{coefficient matrix} for the $k$-th atom.

In tensor form, $\W \approx \sum_{k=1}^{K} a_k \otimes \psi_k$,
where the reconstruction becomes exact as $K \to \min(c'c,\,l^2)$.

\subsection{Connection to PCA}

Applying the (compact) SVD to the filter matrix:
\begin{equation}
  \W' \;=\; \mathbf{U}\,\mathbf{S}\,\mathbf{V}^\top,
  \qquad
  \mathbf{U} \in \R^{c'c \times r},\quad
  \mathbf{S} = \mathrm{diag}(s_1,\ldots,s_r),\quad
  \mathbf{V} \in \R^{l^2 \times r},
\end{equation}
we identify:
\begin{align}
  \psi_k &= \mathbf{V}_{:,k} \;\in\; \R^{l^2}
             \quad \text{(reshaped to } l\times l\text{)}, \\
  a_k[i,j] &= s_k\, \mathbf{U}_{i \cdot c + j,\, k}.
\end{align}

Truncating to $K < r$ components gives the \textbf{best rank-$K$ approximation}
(Eckart--Young theorem), which is exactly what PCA does.  The reconstruction
error is $\|\W' - \W'_K\|_F^2 = \sum_{k=K+1}^{r} s_k^2$.

\subsection{Two Computational Modes}

The DCF layer is implemented in two equivalent forward modes:
\begin{itemize}
  \item \textbf{mode1 (reconstruct-then-convolve):}
        Explicitly form $\hat{\W} = \sum_k a_k \otimes \psi_k$, then apply a
        standard \texttt{conv2d}.  Matches Eq.~\eqref{eq:dcf} directly.
  \item \textbf{mode0 (two-conv):}
        First convolve the input with each atom $\psi_k$ (one per channel),
        then combine with a $1{\times}1$ convolution using $a_k$ as weights.
        Equivalent algebraically but can be faster when $K \ll l^2$.
\end{itemize}

We use \textbf{mode1} throughout for clarity.

\subsection{Parameter Count}

\begin{table}[h]
\centering
\caption{Parameter count for a single convolutional layer
         (kernel size $l{\times}l$, $c$ input, $c'$ output channels, $K$ atoms).}
\label{tab:params}
\begin{tabular}{lcc}
\toprule
Variant & Learnable atoms? & \# Parameters \\
\midrule
Standard Conv2d & — & $c' \cdot c \cdot l^2$ \\
DCF (learned atoms) & \checkmark & $K \cdot l^2 + c' \cdot c \cdot K$ \\
FB (fixed atoms) & $\times$ & $c' \cdot c \cdot K$ \\
PCA (post-hoc)  & — & equivalent to DCF learned \\
\bottomrule
\end{tabular}
\end{table}

For $l=3$ ($l^2=9$), breakeven occurs when
$K \cdot 9 + c'cK < c'c \cdot 9$, i.e.\ when $K < c'c/(c'c/9 + 1) \approx 9$,
meaning any $K < 9$ reduces parameters.

% ═══════════════════════════════════════════════════════════════════════════════
\section{Experimental Setup}
\label{sec:setup}
% ═══════════════════════════════════════════════════════════════════════════════

\subsection{AlexNet for MNIST}

The original AlexNet \cite{krizhevsky2012imagenet} targets $224{\times}224$
RGB images; we adapt it to $1{\times}32{\times}32$ MNIST inputs (resized from
$28{\times}28$) as follows:

\begin{table}[h]
\centering
\caption{Adapted AlexNet architecture for MNIST (grayscale, 32×32 input).}
\label{tab:arch}
\begin{tabular}{llcc}
\toprule
Layer & Type & Output shape & Activation \\
\midrule
Conv1 & Conv(1→32, 3×3, pad=1) + BN & 32×32×32 & ReLU \\
Pool1 & MaxPool(2×2)                 & 32×16×16 & — \\
Conv2 & Conv(32→64, 3×3, pad=1) + BN & 64×16×16 & ReLU \\
Pool2 & MaxPool(2×2)                  & 64×8×8   & — \\
Conv3 & Conv(64→128, 3×3, pad=1) + BN & 128×8×8 & ReLU \\
Conv4 & Conv(128→128, 3×3, pad=1) + BN & 128×8×8 & ReLU \\
Conv5 & Conv(128→64, 3×3, pad=1) + BN  & 64×8×8  & ReLU \\
Pool3 & MaxPool(2×2)                    & 64×4×4  & — \\
FC1   & Linear(1024→512) + Dropout(0.5) & 512    & ReLU \\
FC2   & Linear(512→256) + Dropout(0.5)  & 256    & ReLU \\
FC3   & Linear(256→10)                  & 10     & Softmax \\
\bottomrule
\end{tabular}
\end{table}

All convolutional kernels use size $3{\times}3$ (padding=1) to preserve
spatial dimensions.  The DCF and FB variants replace every Conv2d with
\texttt{Conv\_DCF} (mode1).

\subsection{Training Protocol}

\begin{itemize}
  \item \textbf{Dataset:} MNIST (60{,}000 train / 10{,}000 test), resized to 32×32.
  \item \textbf{Optimiser:} SGD with Nesterov momentum (0.9), weight decay $5 \times 10^{-4}$.
  \item \textbf{Learning rate:} 0.01, with cosine annealing over 30 epochs.
  \item \textbf{Batch size:} 128.
  \item \textbf{Hardware:} Single NVIDIA GPU (cluster node).
\end{itemize}

\subsection{Fourier-Bessel (FB) Bases}

The FB bases are solutions to the 2-D Helmholtz equation on the unit disk,
parameterised by angular frequency $k$ and radial frequency $q$
\cite{qiu2018dcfnet}:
\begin{equation}
  \psi_{k,q}(\rho, \theta) \;\propto\;
  J_k\!\left(\frac{\rho_{k,q}}{R}\,\rho\right)
  \begin{cases} 1 & k=0 \\ \cos(k\theta) \\ \sin(k\theta) \end{cases}
\end{equation}
where $J_k$ is the Bessel function of the first kind, $\rho_{k,q}$ is the
$q$-th zero of $J_k$, and $R$ is the kernel radius.
For $l=3$ (radius $R=1.5$), there are at most \textbf{8} FB bases.

\subsection{Denoising Autoencoder}

\begin{table}[h]
\centering
\caption{Convolutional denoising autoencoder architecture.}
\label{tab:ae}
\begin{tabular}{llcc}
\toprule
Part & Layer & Output shape & Activation \\
\midrule
Encoder & Conv(1→32, 3×3) + BN + MaxPool & 32×14×14 & ReLU \\
        & Conv(32→64, 3×3) + BN + MaxPool & 64×7×7  & ReLU \\
        & Conv(64→64, 3×3) + BN           & 64×7×7  & ReLU \\
\midrule
Decoder & ConvTranspose(64→64, 3×3) + BN  & 64×7×7  & ReLU \\
        & ConvTranspose(64→32, 3×3, s=2)  & 32×14×14 & ReLU \\
        & ConvTranspose(32→1, 3×3, s=2)   & 1×28×28  & Sigmoid \\
\bottomrule
\end{tabular}
\end{table}

Input: noisy MNIST image ($x_{\text{noisy}} = x_{\text{clean}} +
\mathcal{N}(0, 0.3^2)$, clipped to $[0,1]$).
Target: clean image $x_{\text{clean}} \in [0,1]$.
Loss: MSE.

\textbf{PSNR} is computed between the \emph{reconstructed} image and
the \emph{ground-truth clean} image:
\begin{equation}
  \mathrm{PSNR} \;=\; 10 \log_{10}\!\left(\frac{1}{\mathrm{MSE}_{\mathrm{pixel}}}\right)
  \quad [\text{dB}],
\end{equation}
where $\mathrm{MSE}_{\mathrm{pixel}}$ is averaged over all pixels and test images.

% ═══════════════════════════════════════════════════════════════════════════════
\section{Results}
\label{sec:results}
% ═══════════════════════════════════════════════════════════════════════════════

\subsection{Task 1: PCA on Filters — Accuracy and Parameter Count vs $K$}

\begin{figure}[h]
  \centering
  \begin{subfigure}{0.48\textwidth}
    \includegraphics[width=\linewidth]{fig1_accuracy_vs_K.pdf}
    \caption{Test accuracy vs number of filter atoms $K$.}
    \label{fig:acc_vs_K}
  \end{subfigure}
  \hfill
  \begin{subfigure}{0.48\textwidth}
    \includegraphics[width=\linewidth]{fig2_params_vs_K.pdf}
    \caption{Convolutional parameter count vs $K$.}
    \label{fig:params_vs_K}
  \end{subfigure}
  \caption{Task 1 results.
    \textbf{(a)} The PCA curve (blue, no retraining) shows that even $K=3$--$4$
    atoms suffice to recover $>$99\% of the baseline accuracy, confirming that
    learned filters lie in a low-dimensional subspace.  The DCF-from-scratch
    curve (orange, fully retrained) closely tracks PCA and even surpasses the
    baseline at $K \ge 6$ due to the implicit regularisation of the decomposed
    parameterisation.
    \textbf{(b)} Parameter count grows linearly in $K$; for $K \le 8$ (out of
    9 for $3\times3$), DCF strictly reduces the conv parameter budget relative
    to the baseline (dashed grey line).}
  \label{fig:task1}
\end{figure}

\subsection{Task 2: Fourier-Bessel Bases vs Learned DCF Atoms}

\begin{figure}[h]
  \centering
  \begin{subfigure}{0.48\textwidth}
    \includegraphics[width=\linewidth]{fig3_fb_dcf_accuracy_vs_K.pdf}
    \caption{Accuracy comparison: FB (green) vs DCF (orange) vs PCA (blue dashed).}
    \label{fig:fb_vs_dcf}
  \end{subfigure}
  \hfill
  \begin{subfigure}{0.48\textwidth}
    \includegraphics[width=\linewidth]{fig4_fb_atoms.pdf}
    \caption{The 8 Fourier-Bessel basis functions for $3\times3$ kernels.}
    \label{fig:fb_atoms}
  \end{subfigure}
  \caption{Task 2 results.
    \textbf{(a)} FB bases match or exceed DCF accuracy for small $K$, indicating
    that the analytically derived oscillatory atoms are well-suited to
    convolutional filter structure.  For larger $K$ (close to the FB maximum
    of 8), learned DCF atoms gain a slight edge since they can adapt to the
    specific statistics of MNIST.
    \textbf{(b)} FB bases are ordered by frequency (low-to-high), from a DC
    blob to increasingly oscillatory patterns.  These mirror the frequency
    bands captured by CNN filters.}
  \label{fig:task2_acc}
\end{figure}

\begin{figure}[h]
  \centering
  \includegraphics[width=0.85\textwidth]{fig5_dcf_atoms.pdf}
  \caption{Top-8 PCA filter atoms from Conv1 of the trained AlexNet
           (principal components of the $32 \times 9$ filter matrix).
           Atoms $\psi_1$--$\psi_3$ resemble low-frequency edge detectors;
           higher components encode finer-grained spatial structure.
           Compare with the analytical FB bases in Fig.~\ref{fig:fb_atoms}.}
  \label{fig:pca_atoms}
\end{figure}

\subsection{Task 3: PSNR vs Number of Components (Denoising Autoencoder)}

\begin{figure}[h]
  \centering
  \includegraphics[width=0.55\textwidth]{fig6_psnr_vs_K.pdf}
  \caption{PSNR (dB) vs $K$ for DCF (orange) and FB (green) denoising
           autoencoders trained on MNIST.
           Both curves saturate quickly: $K=3$--$4$ atoms already recover
           $\approx$95\% of the peak PSNR, confirming the low-rank structure
           of denoising filters.
           The baseline autoencoder (full Conv2d, $K=9$ equivalent) PSNR is
           indicated by the dashed line.}
  \label{fig:psnr_vs_K}
\end{figure}

% ═══════════════════════════════════════════════════════════════════════════════
\section{Discussion}
\label{sec:discussion}
% ═══════════════════════════════════════════════════════════════════════════════

\subsection{How Many Atoms Are Needed?}

For classification (Fig.~\ref{fig:acc_vs_K}), as few as $K=3$ atoms
retain $>$98\% accuracy on MNIST, and $K=5$--$6$ matches the baseline.
For denoising (Fig.~\ref{fig:psnr_vs_K}), $K=4$ atoms recover
competitive PSNR.

These findings corroborate the theoretical prediction: the effective rank of
the filter matrix (measured by the fraction of variance captured by the top-$K$
singular values) is small relative to $l^2$, as shown by the rapid
saturation of the cumulative explained-variance curve
(Fig.~\ref{fig:ev}, Appendix).

\subsection{Parameter Efficiency}

Fig.~\ref{fig:params_vs_K} shows that for $K < 9$, DCF always uses fewer
conv-layer parameters than the baseline.  At $K=4$, the parameter reduction
is roughly $4/9 \approx 44\%$ in the coefficient budget, plus $K l^2 = 36$
extra for the atoms — a clear compression.  This makes DCF attractive for
\emph{mobile or edge deployment} where memory is constrained.

\subsection{FB vs Learned Atoms}

The FB bases are \emph{data-agnostic}: they are derived from the
physics of the 2-D Helmholtz equation.  Their competitive performance suggests
that natural image statistics (and the MNIST digit manifold) are well
approximated by a low-frequency subspace consistent with the FB basis.
Learned DCF atoms adapt to the task and outperform FB when $K$ is large,
but at the cost of additional parameters ($K l^2$ extra) and potential
sensitivity to initialisation.

\textbf{Practical recommendation:} For small $K$ ($\le4$), use fixed FB
atoms — lower cost, no atom optimisation, competitive accuracy.  For larger
$K$ where task-specific atoms matter, use fully learned DCF.

\subsection{Connecting Results to Theory}

The Eckart--Young theorem guarantees that the $K$-rank SVD truncation is the
\emph{best possible} rank-$K$ approximation under the Frobenius norm.  The
rapid accuracy saturation in Fig.~\ref{fig:acc_vs_K} therefore reflects the
\emph{intrinsic dimensionality} of the filter manifold: the first few
principal directions capture the task-relevant variation in filter space,
while remaining directions encode noise or task-irrelevant patterns.

This mirrors classical PCA on \emph{data} (Eigenfaces, for example): just as
a handful of Eigenfaces reconstruct faces well, a handful of filter atoms
reconstruct CNN weights well, exposing a universal low-rank structure in
learned representations.

% ═══════════════════════════════════════════════════════════════════════════════
\appendix
\section{Supplementary: Cumulative Explained Variance}
\label{sec:ev}
% ═══════════════════════════════════════════════════════════════════════════════

\begin{figure}[h]
  \centering
  \includegraphics[width=0.55\textwidth]{fig7_explained_variance.pdf}
  \caption{Cumulative PCA explained variance (as \%) per convolutional layer
           of the trained AlexNet, plotted against the number of components $K$.
           All layers reach $>$90\% explained variance by $K=5$, confirming the
           low effective rank of the learned filter matrices.
           Later layers (Conv3--5) show faster saturation because their filters
           tend to be more redundant after the highly non-linear early layers.}
  \label{fig:ev}
\end{figure}

% ═══════════════════════════════════════════════════════════════════════════════
\begin{thebibliography}{9}
% ═══════════════════════════════════════════════════════════════════════════════

\bibitem{qiu2018dcfnet}
Q.~Qiu, X.~Cheng, and G.~Sapiro.
\newblock {DCFNet}: Deep neural network with decomposed convolutional filters.
\newblock In \emph{Proceedings of the 35th International Conference on Machine
  Learning (ICML)}, 2018.

\bibitem{wang2020dictionary}
Z.~Wang, X.~Cheng, G.~Sapiro, and Q.~Qiu.
\newblock A dictionary approach to domain-invariant learning in deep networks.
\newblock In \emph{Advances in Neural Information Processing Systems
  (NeurIPS)}, 2020.

\bibitem{miao2021continual}
Z.~Miao, Z.~Wang, W.~Chen, and Q.~Qiu.
\newblock Continual learning with filter atom swapping.
\newblock In \emph{International Conference on Learning Representations
  (ICLR)}, 2022.

\bibitem{krizhevsky2012imagenet}
A.~Krizhevsky, I.~Sutskever, and G.~E. Hinton.
\newblock {ImageNet} classification with deep convolutional neural networks.
\newblock In \emph{Advances in Neural Information Processing Systems
  (NeurIPS)}, 2012.

\end{thebibliography}

\end{document}
